\chapter{Danh sách quy chế đào tạo}
\label{app:quy-che-dao-tao}

Phụ lục này mô tả chi tiết năm văn bản quy chế đào tạo sau đại học được sử dụng làm nguồn dữ liệu cho miền đào tạo trong hệ thống. Các văn bản do Trường Đại học Công nghệ Thông tin -- Đại học Quốc gia Thành phố Hồ Chí Minh ban hành, điều chỉnh toàn bộ hoạt động đào tạo bậc thạc sĩ và tiến sĩ.

% -------------------------------------------------------------------
\section{Quy chế 1: Quy chế đào tạo trình độ thạc sĩ}
\label{sec:app-qc1}

\textbf{Ngày ban hành:} Tháng 9 năm 2021.

\textbf{Cơ quan ban hành:} Hiệu trưởng Trường Đại học Công nghệ Thông tin -- ĐHQG-HCM.

\textbf{Phạm vi áp dụng:} Áp dụng cho toàn bộ chương trình đào tạo thạc sĩ tại Trường Đại học Công nghệ Thông tin, bao gồm các chuyên ngành Khoa học Máy tính, Hệ thống Thông tin, Kỹ thuật Máy tính và các chuyên ngành liên quan.

\textbf{Cấu trúc văn bản:} Văn bản gồm 6 chương với tổng cộng 35 điều, bao quát toàn bộ vòng đời đào tạo từ tuyển sinh đến tốt nghiệp:
\begin{itemize}
  \item Chương I: Những quy định chung (Điều 1--5)
  \item Chương II: Tuyển sinh và nhập học (Điều 6--10)
  \item Chương III: Chương trình và tổ chức đào tạo (Điều 11--18)
  \item Chương IV: Đánh giá kết quả học tập (Điều 19--25)
  \item Chương V: Luận văn thạc sĩ (Điều 26--32)
  \item Chương VI: Điều khoản thi hành (Điều 33--35)
\end{itemize}

\textbf{Nội dung chính:} Quy chế này là văn bản nền tảng điều chỉnh toàn bộ hoạt động đào tạo thạc sĩ. Các nội dung trọng tâm bao gồm: điều kiện dự tuyển và quy trình tuyển sinh; khối lượng học tập tối thiểu và thời gian hoàn thành chương trình; hệ thống tín chỉ, thang điểm và điều kiện công nhận môn học; quy định về cảnh báo học vụ và buộc thôi học; điều kiện đề xuất đề tài và đăng ký bảo vệ luận văn; quy trình thành lập hội đồng chấm luận văn và điều kiện tốt nghiệp.

% -------------------------------------------------------------------
\section{Quy chế 2: Quy định về luận văn thạc sĩ}
\label{sec:app-qc2}

\textbf{Ngày ban hành:} Tháng 3 năm 2022.

\textbf{Cơ quan ban hành:} Hiệu trưởng Trường Đại học Công nghệ Thông tin -- ĐHQG-HCM.

\textbf{Phạm vi áp dụng:} Áp dụng riêng cho giai đoạn thực hiện và bảo vệ luận văn thạc sĩ, bổ sung chi tiết cho các quy định về luận văn tại Chương V của Quy chế đào tạo thạc sĩ.

\textbf{Cấu trúc văn bản:} Văn bản gồm 4 chương với 28 điều:
\begin{itemize}
  \item Chương I: Quy định chung về luận văn (Điều 1--6)
  \item Chương II: Đề xuất và phê duyệt đề tài (Điều 7--12)
  \item Chương III: Thực hiện và đánh giá luận văn (Điều 13--22)
  \item Chương IV: Điều khoản thi hành (Điều 23--28)
\end{itemize}

\textbf{Nội dung chính:} Văn bản quy định chi tiết về: tiêu chuẩn chọn người hướng dẫn và phân công hướng dẫn; yêu cầu nội dung và hình thức trình bày luận văn; thủ tục nộp luận văn và kiểm tra đạo văn; thành phần và thẩm quyền của hội đồng chấm luận văn; quy trình buổi bảo vệ và tính điểm; điều kiện chỉnh sửa luận văn sau bảo vệ và hoàn thiện hồ sơ tốt nghiệp.

% -------------------------------------------------------------------
\section{Quy chế 3: Quy định về học bổng và học phí}
\label{sec:app-qc3}

\textbf{Ngày ban hành:} Tháng 8 năm 2022.

\textbf{Cơ quan ban hành:} Hiệu trưởng Trường Đại học Công nghệ Thông tin -- ĐHQG-HCM, ban hành dựa trên khung học phí của ĐHQG-HCM.

\textbf{Phạm vi áp dụng:} Áp dụng cho tất cả học viên cao học và nghiên cứu sinh đang theo học tại Trường Đại học Công nghệ Thông tin.

\textbf{Cấu trúc văn bản:} Văn bản gồm 3 chương với 20 điều:
\begin{itemize}
  \item Chương I: Học phí và phương thức thanh toán (Điều 1--8)
  \item Chương II: Học bổng và hỗ trợ tài chính (Điều 9--16)
  \item Chương III: Điều khoản thi hành (Điều 17--20)
\end{itemize}

\textbf{Nội dung chính:} Văn bản quy định mức học phí theo tín chỉ và theo học kỳ; thời hạn và hình thức đóng học phí; chính sách hoàn trả học phí khi nghỉ học hoặc bảo lưu; các loại học bổng dành cho học viên xuất sắc, học viên có hoàn cảnh khó khăn và học viên người dân tộc thiểu số; điều kiện và thủ tục đăng ký học bổng; quy định về miễn giảm học phí theo diện chính sách.

% -------------------------------------------------------------------
\section{Quy chế 4: Quy định về nghiên cứu sinh và đào tạo tiến sĩ}
\label{sec:app-qc4}

\textbf{Ngày ban hành:} Tháng 1 năm 2023.

\textbf{Cơ quan ban hành:} Hiệu trưởng Trường Đại học Công nghệ Thông tin -- ĐHQG-HCM, phù hợp với Thông tư 18/2021/TT-BGDĐT của Bộ Giáo dục và Đào tạo.

\textbf{Phạm vi áp dụng:} Áp dụng cho toàn bộ chương trình đào tạo tiến sĩ tại Trường Đại học Công nghệ Thông tin.

\textbf{Cấu trúc văn bản:} Văn bản gồm 7 chương với 42 điều:
\begin{itemize}
  \item Chương I: Những quy định chung (Điều 1--5)
  \item Chương II: Tuyển sinh nghiên cứu sinh (Điều 6--11)
  \item Chương III: Chương trình và kế hoạch đào tạo (Điều 12--18)
  \item Chương IV: Tiểu luận tổng quan và chuyên đề (Điều 19--24)
  \item Chương V: Luận án tiến sĩ (Điều 25--35)
  \item Chương VI: Bảo vệ luận án (Điều 36--40)
  \item Chương VII: Điều khoản thi hành (Điều 41--42)
\end{itemize}

\textbf{Nội dung chính:} Văn bản điều chỉnh toàn bộ quy trình đào tạo tiến sĩ với các nội dung chính: điều kiện và quy trình tuyển sinh nghiên cứu sinh; xây dựng kế hoạch đào tạo cá nhân; yêu cầu về công bố khoa học quốc tế; thực hiện tiểu luận tổng quan tài liệu; quy định về hội đồng chấm luận án cấp cơ sở và cấp trường; điều kiện tốt nghiệp và cấp bằng tiến sĩ.

% -------------------------------------------------------------------
\section{Quy chế 5: Quy định bổ sung về đào tạo sau đại học}
\label{sec:app-qc5}

\textbf{Ngày ban hành:} Tháng 5 năm 2023.

\textbf{Cơ quan ban hành:} Hiệu trưởng Trường Đại học Công nghệ Thông tin -- ĐHQG-HCM.

\textbf{Phạm vi áp dụng:} Bổ sung và cập nhật một số điều khoản của các quy chế trước, áp dụng cho tất cả học viên cao học và nghiên cứu sinh đang theo học kể từ học kỳ 1 năm học 2023--2024.

\textbf{Cấu trúc văn bản:} Văn bản gồm 2 chương với 15 điều bổ sung:
\begin{itemize}
  \item Chương I: Sửa đổi, bổ sung một số điều của các quy chế hiện hành (Điều 1--10)
  \item Chương II: Điều khoản chuyển tiếp và thi hành (Điều 11--15)
\end{itemize}

\textbf{Nội dung chính:} Văn bản này cập nhật một số điều khoản phát sinh từ thực tế triển khai, bao gồm: bổ sung quy định về học phần thay thế và tương đương tín chỉ; cập nhật thủ tục đăng ký học chuyên đề của khóa khác; làm rõ quy định về bảo lưu kết quả học tập và thời gian bảo lưu tối đa; bổ sung hướng dẫn về công nhận tín chỉ từ chương trình nước ngoài; điều chỉnh quy trình xử lý vi phạm học thuật phù hợp với chuẩn mực quốc tế.
